\chapter{绪论} \label{ch:introduction}


\section{研究背景与意义}

预防和治理传染病是一项关乎国家经济和人民健康的重要任务，量化研究疾病传播规律对于制定防控策略至关重要。传染病动力学利用疾病传播规律和外部因素构建数学模型，揭示疾病的演变过程，预测传播规律和趋势，并寻求最佳防控策略，为决策提供依据。流行病仓室模型从机制上解释了特定流行病区别于其他流行病的独特特征，能够很好地理解疾病的传播机制，预测疫情的发展趋势，评估控制措施的有效性。

自新冠肺炎首次出现以来，该疫情席卷全球，导致超过 \num{7.6} 亿例确诊病例和超过 \num{690} 万人死亡~\cite{Data_WHO}，深刻影响了世界的政治经济秩序，改变了普通人的生活。
因此，研究新冠肺炎的动力学模型，以预测其传播趋势和制定最优控制策略，具有重要的理论和实践意义，有助于全球范围内的疫情防控工作。
相比于其他传染病， COVID-19 具有许多独特的特征，如高度传染性、较长的潜伏期、相当比例的无症状患者、不断变异出新毒株、老年人群死亡率高等。
因此，有必要针对 COVID-19 的特征开发具有针对性的数学模型，准确刻画 COVID-19 的动力学性质。
这些结果能帮助政策制定者制定有效的策略，以减缓 COVID-19 的传播。


\section{国内外研究现状}

近年来，国内外研究集中关注以下主要问题：

\begin{itemize}
	\item \textbf{概念综合}~\cite{Age-structured1,Age-structured2,Age-structured3,vaccination1,vaccination2,ZhouBMC2023,Wei2023IDM}：引入年龄结构、疫苗接种、人群流动、非药物干预措施（NPIs）等多种现实因素；
	\item \textbf{大规模数据驱动}~\cite{Shaier2022arXiv,Faith2021medRxiv,Zisad2021Algorithms}：大规模数据追踪人群行为、社交接触和病例报告，从而改进流行病预测和干预策略。
	\item \textbf{人工智能方法}~\cite{Tang2022AAM,Li2021CJE,Zain2021JCSE,Yang2020JTD}：通过人工智能方法分析大规模的流行病数据、改进预测和模拟的准确性，进一步挖掘流行病的隐藏的模式和关联。
\end{itemize}

其中，文献~\cite{Wei2023IDM} 提出了一种SEIRV仓室模型，该模型考虑了年龄结构和疫苗接种情况，特别关注了这些因素与宿主特性以及非药物干预（NPIs）之间复杂的相互作用。
研究基于该模型，深入研究了中国河北省石家庄市农村地区COVID-19的流行病学和动力学特征。

近年来，有学者关注了 LSTM 模型在研究时间序列问题上的应用。
文献~\cite{Huang2018AAAI} 针对临床时间序列数据，采用了基于注意力机制的 LSTM 深度神经网络，摒弃了循环神经网络（RNN）的传统结构。
作者提出了SAnD（Simply Attend and Diagnose）架构，利用掩码自注意机制、位置编码和密集插值策略来处理时间顺序信息。此外，还开发了SAnD的多任务变体，用于同时处理多个诊断任务。
文献~\cite{Xu2020AAAI} 提出了一种 CF-LSTM 框架，用于解决复杂场景下的预测问题。
文献~\cite{Xu2020AAAI} 使用了 LSTM 网络构建一个能够在没有人工专家监督的情况下用于处理用户执行的一系列操作作为数据的模型。

文献~\cite{Zain2021JCSE} 则开发了一种基于时间序列数据的混合 CNN-LSTM 模型，用于COVID-19确诊病例数量的预测。
研究结果表明，将CNN和LSTM模型结合成CNN-LSTM编码器-解码器结构，显著提高了预测性能。
即使在数据量有限的情况下，该模型也能够产生令人满意的预测结果。

此外，文献~\cite{Yang2020JTD} 中将2020年1月23日前后的人口迁徙数据与最新的 COVID-19 流行病学数据整合到 SEIR 模型中。
采用2003年 SARS 病毒的数据训练了一个 LSTM 模型。
结果表明，动态 SEIR 模型在预测 COVID-19 疫情的高峰和规模方面非常有效。

文献~\cite{Delli2022PLOSOne} 结合了神经网络模型和 SEIR 模型创建了一个混合模型。
研究结果表明，混合模型在不同地理层次上改善了 SEIR 模型的性能，可以用于预测ICU的占用率，支持决策制定，例如病人转移和医疗人员以及呼吸机的调配等。

这些研究表明，多领域的探索和创新方法在应对COVID-19等流行病方面发挥着关键作用，为流行病学和公共卫生领域提供了宝贵的见解和工具。


\section{主要工作}

本文的主要工作集中在综合多种方法解决流行病学问题，即将流行病学调查、微分方程理论和神经网络模型相结合，以更好地理解和预测流行病的动态行为。
论文详细探讨了流行病学调查的重要性，包括数据收集、验证模型预测和疫情追踪等方面，以提高流行病模型的准确性和适用性；
微分方程理论在建立流行病模型、参数估计和稳定性分析中的关键作用，以及如何使用微分方程理论来分析模型的平衡点性质和阈值；
最后，论述了神经网络模型在流行病学问题中的应用，包括非线性建模、数据驱动性和时间序列预测等方面，强调神经网络的优势和应用潜力。
综合而言，本文通过整合不同领域的方法，为深入理解和有效应对流行病学问题提供了全面的方法框架和工具。


\section{论文结构}

在引言部分，强调了流行病学研究的重要性和挑战，明确了本研究的目标和主要贡献。
在文献回顾中，深入探讨了流行病学调查方法、传染病建模技术以及神经网络在流行病学中的应用，为后续方法的提出提供了坚实的理论基础。
在方法部分，详细描述了综合方法的设计，包括流行病学调查、微分方程模型和神经网络结合的具体步骤。
随后，通过具体案例研究，将所提出的方法应用于实际流行病事件的分析，展示了方法的实际效果。
结果和讨论部分呈现了案例研究的详细结果，对模型性能进行了评估，并探讨了方法的局限性和未来改进的潜力。